% Part: counterfactuals
% Chapter: minimal-change-semantics
% Section: introduction
%READERNOTE{SOL6-A213: सर्वांत लहान पूर्वांग-गोलकाचा उल्लेख आकृतीपुरता स्पष्ट केला. दोन शेजारच्या दर्शवलेल्या थरांची जवळीक आणि मनमानी गोलकांतील जवळीक वेगळी आहे; सामान्य गोलक-अट सर्वात जवळचे जग नसण्यालाही लागू होते.}

\documentclass[../../../include/open-logic-section]{subfiles}

\begin{document}

\olfileid{cnt}{min}{int}

\olsection{प्रस्तावना}

‘‘ही आगकाडी ओढली असती, तर ती पेटली असती’’
अशा प्रतिवास्तविक सशर्त विधानांचे
स्टालनाकर आणि लुईस यांनी विश्लेषण सुचवले.
अशा वाक्यांच्या सत्यतेच्या अटी
योग्य रीतीने कशा समजाव्यात,
याविषयी त्यांची मांडणी होती.
दोन्ही मांडण्यांमागील कल्पना अशी आहे:
प्रतिवास्तविक सशर्त विधान सत्य आहे का
हे ठरवताना त्याचे पूर्वांग सत्य होण्यासाठी
वास्तविक जगापासून किमान फरक असलेली
संभाव्य जगे विचारात घ्यावी लागतात.
या संभाव्य जगांत उत्तरांग सत्य असेल,
तर प्रतिवास्तविक विधान सत्य असते.
उदाहरणार्थ, माझ्या हातात आगकाडी
आणि आगपेटी आहे असे समजा.
वास्तविक जगात मी फक्त त्यांच्याकडे
पाहतो आणि आगकाडी ओढली तर काय
होईल, याचा विचार करतो.
मी आगकाडी ओढतो अशा पर्यायी जगातला
किमान बदल म्हणजे मी कृती करण्याचे
ठरवतो आणि आगकाडी ओढतो.
हा बदल किमान आहे, कारण आगकाडी ओढणे आणि त्याचे
आवश्यक परिणाम यांखेरीज दुसरे अनावश्यक बदल होत नाहीत: मी त्याच वेळी
हवेत उडी मारत नाही; आगकाडी
ओढल्याने माझ्या केसांनाही आग
लागत नाही; माझ्या बोटांतील सारी
ताकद अचानक नाहीशी होत नाही;
त्याच वेळी कुणी सुपर-सोकरने
माझ्यावर पाणी मारत नाही;
इत्यादी. त्या पर्यायी शक्यतेत
आगकाडी पेटते. म्हणून मी
आगकाडी ओढली असती,
तर ती पेटली असती,
हे सत्य आहे.

या सहज समजणाऱ्या मांडणीला
प्रतिवास्तविक तर्कशास्त्रांची औपचारिक
चिन्हार्थमीमांसा जोडता येते.
प्रतिवास्तविक विधानासाठी लुईस यांनी
‘‘$\boxright$’’ हे चिन्ह, तर
स्टालनाकर यांनी ‘‘$>$’’ हे चिन्ह
वापरले. आपण~$\cif$ वापरू आणि
ते विधानीय तर्कशास्त्रात द्विपदी
संयोजक म्हणून जोडू.
म्हणून $!A \lif !B$ या रूपाच्या
!!{formula}s बरोबरच
$!A \cif !B$ या रूपाची
!!{formula}s ही असतील.
मोडल तर्कशास्त्राच्या संबंधात्मक
चिन्हार्थमीमांसेप्रमाणे, औपचारिक
चिन्हार्थमीमांसा अशा प्रतिमानांवर
आधारलेली असते की ज्यांत
!!{formula}s यांचे मूल्यांकन जगांच्या
संदर्भात केले जाते. त्यात
$\mSat{M}{!A \cif !B}[w]$
याची पूर्ति-अट काही (इतर) जगे~$w'$
यांच्या संदर्भातील
$\mSat{M}{!A}[w']$ आणि
$\mSat{M}{!B}[w']$
यांच्या साहाय्याने दिली जाते.
कोणती $w'$ जगे?
सहज कल्पनेनुसार,
$\mSat{M}{!A}[w']$ सत्य असलेली
आणि~$w$ च्या सर्वाधिक जवळ
असलेली जगे. त्यासाठी
‘‘जवळीक’’ दर्शवणारा
संबंधही प्रतिमानात
समाविष्ट करावा लागतो.

प्रतिवास्तविक परिस्थिती चित्ररूपात
दाखवण्याची एक बोधक पद्धत लुईस
यांनी आणली. प्रत्येक संभाव्य जग
हे इतर जगे समाविष्ट करणाऱ्या
एकमेकांच्या आत असलेल्या गोलकांच्या
प्रणालीच्या केंद्रस्थानी असते;
आपण हे गोलक समकेंद्री
वर्तुळांनी दाखवतो. या आकृतीत शेजारच्या दोन गोलकांच्या मधला थर
समान जवळीक दाखवतो. अधिक सूक्ष्म गोलकांची सामान्य
प्रणाली असताना मनमानी दोन गोलकांमधील सर्व जगे
सारख्याच जवळ असण्याची अट नसते.
आतल्या गोलकात असलेली जगे
अधिक जवळ, तर बाहेरच्या
गोलकात असलेली जगे अधिक
दूर असतात.
\begin{center}
\begin{tikzpicture}[scale=.7]\small
  \spheresystem{5}
  \draw (0,0) node {$w$};
  \propositionintersect{0}{4}{40}{3.5}
  {\tikzset{proposition/.append={smooth,tension=1.4}}
  \proposition[shift={(1.6,-1)}]{90}{1}{30}{4}}
  \path (1.5,3) node[below] {$\formula{B}$};
  \path (3,0) node {$\formula{A}$};
\end{tikzpicture}
\end{center}
या आकृतीतील सर्वाधिक जवळची $!A$-जगे म्हणजे
$!A$ सत्य असलेले किमान एक जग धारण करणाऱ्या
सर्वांत लहान गोलकातील $!A$-जगे~$w'$ (राखाडी भाग).
तो प्रणालीतील सर्वांत लहान गोलकच असावा असे नाही.
असा $!A$-जग धारण करणारा सर्वांत लहान गोलक
सर्व प्रतिमानांत अस्तित्वात असतोच असे नाही;
पुढील विभागातील गोलक-अट अशा प्रकरणांनाही लागू होते.
सहज कल्पनेनुसार,
सर्वाधिक जवळच्या सर्व
$!A$-जगांत $!B$ सत्य
असेल, तर~$w$ येथे
$!A \cif !B$ ची
पूर्ति होते.

\end{document}
